\section{Related Work}\label{sec:related}

\subsection{Zero Trust for Industrial and OT Systems}
Zero Trust architecture is by now well surveyed for IoT. Azad et al.\ cover continuous verification and micro-segmentation \cite{azad2024verify}, James et al.\ review authentication and authorization schemes specifically under a Zero Trust model \cite{james2024authz}, and Nahar et al.\ examine the 6G setting, where device density sharpens per-request overhead concerns \cite{nahar2024survey}. Fern\'andez and Brazhuk caution that ``Zero Trust'' is sometimes a label without a falsifiable architecture behind it \cite{fernandez2024critical}, and Phiayura and Teerakanok supply a migration framework with a maturity model \cite{phiayura2023comprehensive}. For OT in particular, Song et al.\ propose a Zero Trust architecture for critical operational technology systems \cite{song2025ztot}, Zanasi et al.\ a flexible architecture for IIoT \cite{zanasi2024flexible}, and Kumar and Sharma integrate blockchain with dynamic access control for Industry~4.0 without a DID-anchored registry \cite{kumar2024industry40}. These works establish the principles; they leave open who owns revocation in a multi-party plant and how a decision should be audited afterwards.

\subsection{Blockchain Identity and Authentication for IIoT}
Surveys of blockchain-based decentralized identification note limits around revocation latency and interoperability \cite{hosseini2023blockchain}, and a broad review of blockchain in IIoT discusses security and privacy opportunities alongside scalability and integration costs \cite{essaid2025iiotbc}. Aguru and Erukala propose a blockchain self-sovereign-identity framework for standardized IoT identification \cite{aguru2025ssi}, Folino et al.\ tie identity to a physically unclonable function as a hardware root of trust \cite{folino2025puf}, and Ren et al.\ combine self-sovereign identity with Zero Trust for AIoT using selective-disclosure credentials \cite{ren2025aiot}. On the industrial authentication side, Zhai et al.\ design a lightweight scheme for IIoT that combines blockchain with Chebyshev chaotic maps to suit constrained devices \cite{zhai2025chebyshev}. DAXiot shows that DID challenge--response for MQTT can run without a blockchain at about 21\% connection overhead over mutual TLS, but then has no shared revocation ledger \cite{philipp2023daxiot}. For revocation itself, Mazzocca et al.\ propose EVOKE, a cryptographic-accumulator mechanism with about 1.5\,KB of per-device state that spreads revocation epidemically among devices \cite{mazzocca2024evoke}; it complements an on-chain status check rather than replacing it. Kumar et al.\ run a privacy-preserving identity system on a private Ethereum proof-of-authority (PoA) chain in which a decentralized voter group governs identity updates and revocations, and report roughly 3--4\,s for registration and under 5\,s for revocation on the live network \cite{kumar2026privacy}; it is the closest published precedent for our governance contract and the only live-chain figure we compare against.

\subsection{Blockchain--Zero Trust Combinations}
Nie et al.\ pair smart-contract policy enforcement with inner-product encryption for fine-grained access in 6G IoT \cite{nie2025zeroTrust}. Salim et al.\ propose ZT-BlocIoT with dynamic trusted gateways and reinforcement-learning-driven sharding, reporting 2{,}500 transactions per second on 700 nodes at 42\,ms latency \cite{salim2026ztblociot}. Wang et al.\ observe that even Zero Trust--inspired schemes for medical IoT typically authenticate once and then trust the session, and propose blockchain-backed continuous authentication \cite{wang2026medicaliot}; the same gap is what motivates our per-request design in industry. Mukta et al.\ delegate Zero Trust access control through a blockchain, which is close to our policy layer but concerns delegation rather than the device identity lifecycle \cite{mukta2026ztdelegation}. For smart factories, Farhan et al.\ describe a conceptual Zero Trust framework on a consortium chain \cite{farhan2026ztfsf}, and Potturi et al.\ combine AI risk scoring with blockchain-anchored federated identity across cloud, edge and IoT \cite{potturi2026hybrid}. Jamwal et al.\ conclude that blockchain-enabled Zero Trust is still mostly proposed at survey level rather than delivered as benchmarked open code \cite{jamwal2026survey}, and a bibliometric analysis finds blockchain with Zero Trust identity management to be one of the field's most active clusters \cite{liu2024dissecting}. Usama et al.\ add decentralized identity, zero-knowledge proofs and Byzantine-resilient federated learning against AI-driven identity attacks \cite{usama2026multilayer}. Federated learning is also studied as a way to share threat intelligence across industrial sites without exposing raw data \cite{tom2025flcti}; we do not use it, but it is a plausible source for the anomaly signal in Section~\ref{sec:signals}.

\subsection{Off-Chain Dependencies, Cost and Oracles}
Because on-chain controls rest on off-chain facts, the reliability of the data that reaches a contract matters. Shujaa et al.\ stress that device attestation, secure boot and key custody remain off-chain dependencies of any blockchain-IoT security design \cite{shujaa2025enhancing}, and Hassan et al.\ survey the oracle problem, that is, how to bring trustworthy external data to a contract \cite{hassan2023oracle}; our trust oracle is an instance of it, and Section~\ref{sec:disc} discusses its single-key weakness. Smart-contract defects are another risk, reviewed by Wei et al.\ \cite{wei2024quality}. On cost, gas has been studied for Ethereum smart contracts \cite{ta2024gascost} and for IoT access control \cite{hasan2023smartcontract}, DPoS and PBFT-based consensus choices have been proposed for IoT and clinical settings \cite{haque2024scalable,peepliwal2024ztab}, and gasless relaying can spare constrained devices from paying fees \cite{madhwal2025gasless}. Ren et al.\ show for vehicular networks that published trust values can themselves link pseudonyms and propose differential privacy to prevent it \cite{ren2026ditdp}, which matters for any design that exposes scores on a shared ledger.

\subsection{Positioning}
Table~\ref{tab:relatedwork} compares the most relevant systems on five properties that matter in a plant: on-chain revocation, continuous per-request authorization, consequence-aware (tiered) policy, handling of legacy and recall scenarios, and a released artifact. The assessments rest on the abstracts and the parts of the full texts we could read; ``not stated'' means we found no statement, not that the property is absent. To our knowledge none of the reviewed systems reports all five.

\begin{table*}[t]
\centering
\caption{Comparison of representative blockchain and Zero Trust systems for (I)IoT identity}
\label{tab:relatedwork}
\scriptsize
\begin{tabular}{p{3.2cm}p{3.6cm}p{1.6cm}p{1.6cm}p{1.6cm}p{2.1cm}p{1.6cm}}
\toprule
\textbf{System} & \textbf{Core mechanism} & \textbf{On-chain revocation} & \textbf{Per-request authZ} & \textbf{Tiered policy} & \textbf{Legacy / recall} & \textbf{Open artifact} \\
\midrule
Zhai et al.\ 2025 \cite{zhai2025chebyshev} & Blockchain + Chebyshev-map authentication (IIoT) & Not stated & Not stated & Not stated & Lightweight devices & Not stated \\
Song et al.\ 2025 \cite{song2025ztot} & ZTA for critical OT & Not stated & Not stated & Not stated & Not stated & Not stated \\
Zanasi et al.\ 2024 \cite{zanasi2024flexible} & Flexible IIoT ZTA & No & Yes & Partial & Not stated & Not stated \\
Kumar \& Sharma 2024 \cite{kumar2024industry40} & Blockchain + dynamic access control (Industry 4.0) & Not stated & Yes & Not stated & Not stated & Not stated \\
Farhan et al.\ 2026 \cite{farhan2026ztfsf} & ZT on consortium chain (smart factory) & Conceptual & Conceptual & Not stated & Not stated & Not stated \\
Potturi et al.\ 2026 \cite{potturi2026hybrid} & AI risk score + chain-anchored identity & Yes & Yes & Not stated & Not stated & Not stated \\
Salim et al.\ 2026 \cite{salim2026ztblociot} & ZT-BlocIoT (trust + DRL sharding) & Yes & Yes & Not stated & Not stated & Not stated \\
Nie et al.\ 2025 \cite{nie2025zeroTrust} & ZT + inner-product encryption (6G) & Yes & Yes & Fine-grained attributes & Not stated & Not stated \\
Wang et al.\ 2026 \cite{wang2026medicaliot} & Blockchain continuous auth.\ (medical IoT) & Yes & Yes & No & Not stated & Not stated \\
Kumar et al.\ 2026 \cite{kumar2026privacy} & Private PoA identity + voter group & Yes (voter group) & Not stated & No & Not stated & Not stated \\
\textbf{This work} & \textbf{DID registry, tiered ZT contracts, quorum governance, hybrid enforcement} & \textbf{Yes (quorum-governed, batch)} & \textbf{Yes (every call)} & \textbf{Yes (4 tiers, gate, strict)} & \textbf{Proxy identity; batch recall; outage mode} & \textbf{Yes (Solidity + Python)} \\
\bottomrule
\end{tabular}
\end{table*}
